\uuid{jTB5}
\titre{Constante de normalisation d'une densité}
\auteur{Maxime Nguyen}
\organisation{AMSCC}
\datecreate{2026-11-24}
\niveau{L2}
\module{Probabilité}
\chapitre{Probabilité continue}
\sousChapitre{Densité de probabilité}
\difficulte{1}
\theme{densité, constante de normalisation}

\contenu{
\begin{question}
\qtype{onlyone}
On note $\mathbf{1}_{[0;1]}$ l'indicatrice de $[0;1]$ : elle vaut $1$ sur $[0;1]$ et $0$ ailleurs.
Pour quelle valeur de la constante $c > 0$ la fonction $f(x) = c\,x\,\mathbf{1}_{[0;1]}(x)$ est-elle une densité de probabilité ?
\begin{answers}
\bad{$c = 1$}
\bad{$c = \frac12$}
\good{$c = 2$}
\bad{Aucune : la fonction atteint son maximum $c$ en $x = 1$, et il faudrait à la fois $c \leq 1$ et $\int f = 1$.}
\end{answers}
\begin{explanations}
$\int_0^1 c\,x\,\mathrm{d}x = c\left[\frac{x^2}{2}\right]_0^1 = \frac{c}{2} = 1$,
donc $c = 2$, et $f$ vaut $2$ en $x = 1$ : c'est permis.

A : c'est la plus grande valeur de $c$ pour laquelle $f \leq 1$, ou le résultat d'une primitive fausse ($x^2$ au lieu de $\frac{x^2}{2}$).

B : l'équation $\frac{c}{2} = 1$ mal résolue.

D : la condition $c \leq 1$ n'existe pas ; une densité peut dépasser $1$.

Contre-sens visé : croire qu'une densité est majorée par 1, et choisir $c$
pour respecter cette fausse contrainte au lieu d'écrire $\int f = 1$.
\end{explanations}
\end{question}
}
